\BourbakiExerciseGroup

\begin{enumerate}
\def\labelenumi{\arabic{enumi})}
\tightlist
\item
  Let \(K_{0}\) be a field of characteristic p \textgreater{} 0, \(\mathbf{K} = \mathbf{K}_{0}(\mathbf{X}, \mathbf{Y})\) the field of rational fractions in two indeterminates over \(K_{0}\) , U and V two indeterminates, \(\Omega\) an algebraically closed extension of \(\mathbf{K}(\mathbf{U}, \mathbf{V})\) , \(\mathbf{E} = \mathbf{K}(\mathbf{U}, u)\) , where \(u^{p} = X + YU^{p}\) and \(F = K(V, v)\) , where \(v^{p} = X + YV^{p}\) ; E and F are not separable over K (V, p.178, Ex. 6, c)).
\end{enumerate}

\begin{enumerate}
\def\labelenumi{\alph{enumi})}
\item
  Show that \(\mathbf{K}\) is relatively algebraically closed in \(\mathbf{E}\) and in \(\mathbf{F}\) (if \(x \in \mathbf{E}\) is algebraic over \(\mathbf{K}\) , show that we must have \(x^p \in \mathbf{K}\) ; if \(x \notin \mathbf{K}\) , then \(\mathbf{E} = \mathbf{K}(\mathbf{U}, x)\) , deduce that then \(\mathbf{X}^{1/p}\) and \(\mathbf{Y}^{1/p}\) would lie in \(\mathbf{K}(x)\) ).
\item
  Show that E and F are linearly disjoint over K but that K is not relatively algebraically closed in \(\mathbf{K}(\mathbf{E} \cup \mathbf{F})\) . (Prove that \(X^{1/p} \in \mathbf{K}(\mathbf{E} \cup \mathbf{F})\) ; deduce that v cannot belong to \(\mathbf{E}(\mathbf{V})\) , and conclude from this that E and F are linearly disjoint over K.)
\end{enumerate}

\begin{enumerate}
\def\labelenumi{\arabic{enumi})}
\setcounter{enumi}{1}
\tightlist
\item
  Let \(\mathbf{K}\) be a field of characteristic \(p > 0\) , \(\Omega\) an extension of \(\mathbf{K}\) and \(\mathbf{E}\) , \(\mathbf{F}\) two subextensions of \(\Omega\) , algebraically disjoint over \(\mathbf{K}\) .
\end{enumerate}

\begin{enumerate}
\def\labelenumi{\alph{enumi})}
\tightlist
\item
  Let \(F_{s}\) be the relative separable closure of K in F. Show that \(E(F_{s})\) is the relative separable closure of E in \(E(F)\) . (Reduce to the case where \(F_{s}=K\) ; let B be a transcendence basis of E over K and \(K'\) the relative algebraic closure of K in F, which is radical over F; then (V, p.~137) \(K'(B)\) is the relative algebraic closure of \(K(B)\) in \(F(B)\) and is p-radical over \(K(B)\) . If \(x\in E(F)\) is algebraic over E, there exists an integer \(r\geqslant0\) such that \(x^{p'}\in M\) , where M is a separable extension of finite degree of \(F(B)\) having a basis \((u_{i})_{1\leqslant i\leqslant n}\) over \(\mathrm{F(B)}\) consisting of elements of E. If \(y = \sum_{i}b_{i}u_{i}\) with \(b_{i}\in \mathrm{F(B)}\) , show
\end{enumerate}

by considering the F(B)-isomorphisms of M into an algebraic closure of E(F) that the \(b_{i}\) belong to K'(B) and deduce that y is p-radical over E.)

\begin{enumerate}
\def\labelenumi{\alph{enumi})}
\setcounter{enumi}{1}
\tightlist
\item
  Suppose that E and F are linearly disjoint over K, that K is relatively algebraically closed in F and that E is a separable extension of K. Show that E is relatively algebraically closed in E(F). (Reduce to the case where E is a separable algebraic extension of K, with the help of V, p.~138. Deduce from a) that the relative algebraic closure of E in E(F) is p-radical over E and then use the fact that if \((a_{\lambda})\) is a basis of E over K, the \(a_{\lambda}^{p'}\) also form a basis of E over K.)
\end{enumerate}

\begin{enumerate}
\def\labelenumi{\arabic{enumi})}
\setcounter{enumi}{2}
\tightlist
\item
  Let E be a field, G a group of automorphisms of E and K the field of invariants of G. a) Show that if there exists no normal subgroup of finite index in G other than G, then E is a regular extension of K.
\end{enumerate}

\begin{enumerate}
\def\labelenumi{\alph{enumi})}
\setcounter{enumi}{1}
\tightlist
\item
  Deduce from a) that if G is isomorphic to a finite product of infinite simple groups, then E is a regular extension of K.
\end{enumerate}

\begin{enumerate}
\def\labelenumi{\arabic{enumi})}
\setcounter{enumi}{3}
\tightlist
\item
  Let K be a field and E, F two extensions of K.
\end{enumerate}

\begin{enumerate}
\def\labelenumi{\alph{enumi})}
\item
  Suppose that K is relatively algebraically closed in F. Show that if \((L, u, v)\) , \((L', u', v')\) are two composite extensions of E and F such that \(u(E)\) and \(v(F)\) (resp \(u'(E)\) and \(v'(F)\) ) are algebraically disjoint over K, then these two composite extensions are isomorphic. (Consider first the case where E is a pure transcendental extension of K; then use Ex. 2, a), as well as Ex. 6, b) of V, p.~154)
\item
  Let \(E_{s}\) , \(F_{s}\) be the relative separable closures of K in E and F respectively. Suppose that all the composite extensions of \(E_{s}\) and \(F_{s}\) are isomorphic. Show that if \((L, u, v)\) and \((L', u', v')\) are two composite extensions of E and F such that \(u(E)\) and \(v(F)\) (resp. \(u'(E)\) and \(v'(F)\) ) are algebraically disjoint over K, then these two composite extensions are isomorphic (use a) and Ex. 11 of V, p.~173).
\end{enumerate}

\begin{enumerate}
\def\labelenumi{\arabic{enumi})}
\setcounter{enumi}{4}
\tightlist
\item
  Let K be a field and p its characteristic exponent. For every integer n \textgreater{} 0 and prime to p, denote by \(\varphi_{\mathrm{K}}(n)\) the degree over K of \(\mathrm{R}_{n}(\mathrm{K}) = \mathrm{K}(\mu_{n})\) , where \(\mu_{n}\) denotes the group of n-th roots of unity in an algebraic closure of K. For example \(\varphi_{\mathrm{Q}}(n) = \varphi(n)\) (V, p.~84, Th. 2); if K is finite with q elements, \(\varphi_{\mathrm{K}}(n)\) is the order of the residue class of q in \((\mathbf{Z}/n\mathbf{Z})^{*}\) (V, p.~97, Cor.); if K is algebraically closed, \(\varphi_{\mathrm{K}}(n) = 1\) for all n.
\end{enumerate}

\begin{enumerate}
\def\labelenumi{\alph{enumi})}
\item
  Let E be the relative algebraic closure of the prime field \(K_{0}\) in K. Show that \(\varphi_{\mathrm{K}}(n) = \varphi_{\mathrm{E}}(n)\) .
\item
  Suppose that K is a finitely generated extension of \(K_{0}\) . Let m be an integer \textgreater0; then there exists an integer N \textgreater0 divisible by every n such that \(\varphi_{\mathrm{K}}(n) \leqslant m\) .
\item
  Let \(s\) be an element of \(\mathrm{GL}_m(\mathbf{K})\) of finite order. Show that \(s^N\) has order a power of \(p\) . (We have \(s = tu = ut\) , where \(u\) has order a power of \(p\) and \(t\) has order \(n\) prime to \(p\) . The element \(t\) is root of a separable polynomial \(\mathbf{X}^n - 1\) as well as of its characteristic polynomial; deduce that the subalgebra \(\mathbf{A}\) of \(\mathbf{M}_m(\mathbf{K})\) generated by \(t\) is etale, of degree \(\leqslant m\) , isomorphic to a product of extensions \(\mathbf{R}_d(\mathbf{K})\) , where \(d\) runs over certain divisors of \(n\) including \(n\) itself. Therefore \(\varphi_{\mathbf{K}}(n) \leqslant m\) ; use \(b\) ).)
\end{enumerate}
% semantic-fragments: legacy-input-seam
\relax{}
